\documentclass[11pt, a4paper]{article}

% --- UNIVERSAL PREAMBLE BLOCK ---
\usepackage[a4paper, top=2.5cm, bottom=2.5cm, left=2cm, right=2cm]{geometry}
\usepackage{fontspec}

\usepackage[portuguese, bidi=basic, provide=*]{babel}

\babelprovide[import, onchar=ids fonts]{portuguese}
\babelprovide[import, onchar=ids fonts]{english}

% Set default/Latin font to Sans Serif in the main (rm) slot
\babelfont{rm}{Noto Sans}
% Assign a specific font for Portuguese text
\babelfont[portuguese]{rm}{Noto Sans}

% Add because main language is not English
\usepackage{enumitem}
\setlist[itemize]{label=-}

% --- ADDITIONAL PACKAGES ---
\usepackage{booktabs}
\usepackage{tabularx}

\begin{document}

\begin{center}
    \Large \textbf{Análise de Variáveis Quantitativas para Tomada de Decisão} \\
    \vspace{0.2cm}
    \large \textit{Estudo de Caso: Controladoria Digital}
\end{center}
\vspace{1cm}

\renewcommand{\arraystretch}{1.8} % Aumenta o espaçamento entre as linhas da tabela
\noindent
\begin{tabularx}{\textwidth}{@{} l >{\raggedright\arraybackslash}X >{\raggedright\arraybackslash}X @{}}
\toprule
\textbf{Cenário} & \textbf{Variáveis quantitativas} & \textbf{Como essa variável auxilia a tomada de decisão?} \\
\midrule
\textbf{Operacional} & 
Horas trabalhadas por projeto; Taxa de cumprimento de prazos (SLA de atendimento). & 
Guia as lideranças na distribuição diária das tarefas. Permite identificar gargalos na rotina da equipa, realocar analistas sobrecarregados e garantir que as obrigações diárias dos clientes sejam entregues a tempo e sem ocorrência de multas. \\
\midrule
\textbf{Tático} & 
Margem de lucro segmentada por tipo de serviço; Índice de inadimplência (percentagem de recebíveis em atraso). & 
Permite aos coordenadores avaliar o que está a funcionar a médio prazo. Ajuda a reestruturar réguas de cobrança, a otimizar processos intermédios de entrega e a direcionar os esforços comerciais para os serviços que apresentam melhor rentabilidade e menor atrito. \\
\midrule
\textbf{Estratégico} & 
Custo de Aquisição de Clientes (CAC); Lifetime Value (LTV); Taxa de evasão da carteira (Churn rate). & 
Fundamenta as decisões de longo prazo da direção sobre o rumo do negócio. Indica, de forma clara, se a empresa deve investir em novas tecnologias, alterar a estrutura global de preços ou focar a sua expansão em novos nichos de mercado. \\
\bottomrule
\end{tabularx}

\end{document}